\documentclass[11pt, a4paper]{article}
\usepackage[a4paper, top=2.5cm, bottom=2.5cm, left=2cm, right=2cm]{geometry}
\usepackage{fontspec}
\usepackage[english, bidi=basic, provide=*]{babel}
\babelprovide[import, onchar=ids fonts]{english}
\babelfont{rm}{TeX Gyre Termes}
\babelfont{sf}{TeX Gyre Heros}
\babelfont{tt}{TeX Gyre Cursor}
\usepackage{enumitem}
\usepackage{amsmath}

\title{\textbf{Comprehensive Exam: Chapters 1--5 \\ 100 Advanced Multiple Choice Questions}}
\author{}
\date{}

\begin{document}

\maketitle

\begin{enumerate}

% ---------------------------------------
% CHAPTER 1: Biology & Scientific Inquiry (1-20)
% ---------------------------------------

\item \textbf{Which of the following represents an emergent property that cannot be observed at the preceding level of biological organization?}
\begin{enumerate}[label=\Alph*.]
    \item A single heart muscle cell contracting in a petri dish.
    \item The amino acid sequence of a transport protein.
    \item The ability of an intact chloroplast to perform photosynthesis, whereas a mixture of its isolated molecules cannot.
    \item The replication of DNA inside a eukaryotic nucleus.
    \item The ratio of carbon to hydrogen in a glucose molecule.
\end{enumerate}
\textbf{Answer:} C \\
\textbf{Explanation:} Emergent properties arise from the specific arrangement and interactions of parts as complexity increases, which is why an intact chloroplast can photosynthesize but its disorganized molecules cannot.

\item \textbf{A researcher studies a complex metabolic pathway by disabling one enzyme at a time and observing the effects on the individual cell. This approach is best described as:}
\begin{enumerate}[label=\Alph*.]
    \item Systems biology
    \item Reductionism
    \item Inductive reasoning
    \item Positive feedback
    \item Genomics
\end{enumerate}
\textbf{Answer:} B \\
\textbf{Explanation:} Breaking down a complex system to study its simpler components individually is the definition of reductionism.

\item \textbf{In a desert ecosystem, kangaroo rats conserve water by producing highly concentrated urine. This physiological adaptation is ultimately the result of:}
\begin{enumerate}[label=\Alph*.]
    \item Goal-directed evolution toward perfection.
    \item Conscious physiological adjustment by the rats.
    \item The rats mutating their genes due to the dry environment.
    \item Unequal reproductive success over generations among individuals with varying genetic traits.
    \item The flow of energy from the sun to the rats.
\end{enumerate}
\textbf{Answer:} D \\
\textbf{Explanation:} Natural selection operates via unequal reproductive success; individuals with heritable traits suited to the environment produce more offspring.

\item \textbf{Blood clotting is an example of positive feedback because:}
\begin{enumerate}[label=\Alph*.]
    \item The end product speeds up its own production.
    \item It is beneficial to the organism's survival.
    \item The response reduces the initial stimulus.
    \item It relies on a continuous input of energy.
    \item It is regulated exclusively by gene expression.
\end{enumerate}
\textbf{Answer:} A \\
\textbf{Explanation:} In positive feedback, the response amplifies the stimulus, accelerating the process until completion (like platelets attracting more platelets).

\item \textbf{Which of the following domains contains organisms that possess membrane-bound organelles?}
\begin{enumerate}[label=\Alph*.]
    \item Bacteria only
    \item Archaea only
    \item Eukarya only
    \item Both Bacteria and Archaea
    \item Both Archaea and Eukarya
\end{enumerate}
\textbf{Answer:} C \\
\textbf{Explanation:} Eukarya is the only domain containing eukaryotic cells, which are defined by having membrane-bound organelles and a true nucleus.

\item \textbf{Darwin's theory of natural selection heavily relied on which of the following observable facts?}
\begin{enumerate}[label=\Alph*.]
    \item The universality of the DNA genetic code.
    \item The mechanism of gene transcription.
    \item The ability of populations to produce more offspring than the environment can support.
    \item The exact age of the Earth as determined by radiometric dating.
    \item The spontaneous generation of simple cells.
\end{enumerate}
\textbf{Answer:} C \\
\textbf{Explanation:} Overpopulation leading to competition was a crucial observation Darwin used to infer unequal reproductive success.

\item \textbf{A scientist formulates the following statement: "If nitrogen is a limiting nutrient for grass growth, then adding nitrogen fertilizer will increase grass mass." This is an example of:}
\begin{enumerate}[label=\Alph*.]
    \item Inductive reasoning
    \item Deductive reasoning
    \item A scientific theory
    \item Qualitative data
    \item Emergent logic
\end{enumerate}
\textbf{Answer:} B \\
\textbf{Explanation:} Deductive reasoning flows from a general premise to a specific, testable prediction, often using "If... then" logic.

\item \textbf{Why is a control group critical to a valid scientific experiment?}
\begin{enumerate}[label=\Alph*.]
    \item It proves that the hypothesis is absolutely true.
    \item It provides a baseline for comparison by isolating the effect of the independent variable.
    \item It eliminates all variables from the experimental environment.
    \item It guarantees that the data collected will be quantitative.
    \item It allows the researcher to test multiple hypotheses simultaneously.
\end{enumerate}
\textbf{Answer:} B \\
\textbf{Explanation:} A control group differs from the experimental group by only the factor being tested, allowing researchers to cancel out the effects of unwanted environmental variables.

\item \textbf{Which of the following is an example of a strictly qualitative observation?}
\begin{enumerate}[label=\Alph*.]
    \item The plant grew 5 centimeters in one week.
    \item The solution turned from clear to bright blue.
    \item The reaction produced 15 grams of precipitate.
    \item The water boiled at $98^\circ\text{C}$.
    \item The population size increased by 15\%.
\end{enumerate}
\textbf{Answer:} B \\
\textbf{Explanation:} Qualitative data consist of descriptions rather than numerical measurements.

\item \textbf{Which statement best distinguishes a scientific theory from a hypothesis?}
\begin{enumerate}[label=\Alph*.]
    \item A theory is an educated guess, while a hypothesis is a proven fact.
    \item A theory can never be disproven, while a hypothesis is easily falsifiable.
    \item A theory is broad in scope, supported by extensive evidence, and generates new hypotheses.
    \item A hypothesis applies to large populations, while a theory applies only to single organisms.
    \item There is no difference; the terms are scientifically interchangeable.
\end{enumerate}
\textbf{Answer:} C \\
\textbf{Explanation:} Theories are comprehensive, highly supported explanatory frameworks, whereas hypotheses are narrow, testable predictions.

\item \textbf{The flow of genetic information inside a cell fundamentally relies on:}
\begin{enumerate}[label=\Alph*.]
    \item Converting proteins directly into DNA.
    \item The exact replication of lipid bilayers.
    \item Transcribing DNA into RNA, which is then translated into protein.
    \item Reversing the sequence of amino acids to build RNA.
    \item Positive feedback loops in the mitochondria.
\end{enumerate}
\textbf{Answer:} C \\
\textbf{Explanation:} Gene expression universally follows the central dogma: DNA $\rightarrow$ RNA $\rightarrow$ Protein.

\item \textbf{An ecologist tracking the flow of energy through a forest ecosystem will observe that energy:}
\begin{enumerate}[label=\Alph*.]
    \item Cycles infinitely between plants and decomposers.
    \item Is completely converted into matter by apex predators.
    \item Enters as light and eventually exits as heat.
    \item Is generated spontaneously by producers.
    \item Is strictly conserved without any loss.
\end{enumerate}
\textbf{Answer:} C \\
\textbf{Explanation:} Unlike chemical nutrients which cycle, energy flows through an ecosystem in one direction, primarily entering as light and exiting as heat.

\item \textbf{Which of the following hypotheses is inherently outside the realm of scientific inquiry?}
\begin{enumerate}[label=\Alph*.]
    \item High dietary sugar increases the risk of heart disease.
    \item Invisible, undetectable energy fields guide bird migration.
    \item A specific gene mutation causes antibiotic resistance in bacteria.
    \item Increasing ocean temperatures cause coral bleaching.
    \item Dinosaurs were driven to extinction by an asteroid impact.
\end{enumerate}
\textbf{Answer:} B \\
\textbf{Explanation:} Hypotheses involving supernatural or explicitly undetectable phenomena cannot be tested or falsified, placing them outside science.

\item \textbf{The unity of life is best illustrated by:}
\begin{enumerate}[label=\Alph*.]
    \item The vast number of different insect species.
    \item The diverse beak shapes of Galápagos finches.
    \item The universal use of DNA as the genetic material.
    \item The presence of both aquatic and terrestrial ecosystems.
    \item The competition between predators and prey.
\end{enumerate}
\textbf{Answer:} C \\
\textbf{Explanation:} Unity in biology refers to shared characteristics derived from common ancestry, such as the universal genetic code.

\item \textbf{A biology student claims that evolutionary adaptation implies organisms intentionally modify their traits to survive. What is the fundamental flaw in this reasoning?}
\begin{enumerate}[label=\Alph*.]
    \item Organisms only modify their traits to attract mates.
    \item Evolution occurs via the selection of naturally occurring, heritable variations, not by intentional effort.
    \item Adaptation only occurs at the cellular level, not the organismal level.
    \item The environment forces mutations to happen exactly when needed.
    \item Traits are never heritable.
\end{enumerate}
\textbf{Answer:} B \\
\textbf{Explanation:} Natural selection is a passive filtering process of existing genetic variation, not a conscious, goal-directed effort by the organism.

\item \textbf{Which domain of life contains organisms known for inhabiting extreme environments, such as boiling hot springs and highly saline lakes?}
\begin{enumerate}[label=\Alph*.]
    \item Fungi
    \item Bacteria
    \item Protista
    \item Archaea
    \item Animalia
\end{enumerate}
\textbf{Answer:} D \\
\textbf{Explanation:} Domain Archaea includes many extremophiles that thrive in harsh environments.

\item \textbf{What is the primary difference between science and technology?}
\begin{enumerate}[label=\Alph*.]
    \item Science is theoretical; technology is hypothetical.
    \item Science aims to apply knowledge to solve problems; technology aims to understand nature.
    \item Science aims to understand natural phenomena; technology aims to apply that knowledge for a specific purpose.
    \item Science requires peer review; technology does not.
    \item Science only uses deductive reasoning.
\end{enumerate}
\textbf{Answer:} C \\
\textbf{Explanation:} Science is the pursuit of knowledge for its own sake, while technology applies scientific discoveries to create tools and solutions.

\item \textbf{Bioinformatics is a vital tool for systems biology because it:}
\begin{enumerate}[label=\Alph*.]
    \item Replaces the need for microscopes.
    \item Synthesizes new elements in the laboratory.
    \item Uses computational power to analyze massive datasets, like entire genomes.
    \item Proves hypotheses using pure mathematics.
    \item Eliminates experimental error.
\end{enumerate}
\textbf{Answer:} C \\
\textbf{Explanation:} The vast amount of data generated by sequencing whole genomes (genomics) requires sophisticated software to organize and analyze.

\item \textbf{In the hierarchy of biological organization, which level is immediately above the "Community" level?}
\begin{enumerate}[label=\Alph*.]
    \item Biosphere
    \item Ecosystem
    \item Population
    \item Organism
    \item Tissue
\end{enumerate}
\textbf{Answer:} B \\
\textbf{Explanation:} The hierarchy moves from Community (all living species in an area) to Ecosystem (all living species plus the nonliving environment).

\item \textbf{The concept of "descent with modification" directly accounts for:}
\begin{enumerate}[label=\Alph*.]
    \item Only the diversity of life.
    \item Only the unity of life.
    \item Both the unity and diversity of life.
    \item The spontaneous generation of life.
    \item The equal sharing of electrons in covalent bonds.
\end{enumerate}
\textbf{Answer:} C \\
\textbf{Explanation:} Descent implies unity (shared ancestry), while modification implies diversity (accumulation of adaptations).

% ---------------------------------------
% CHAPTER 2: Chemistry (21-40)
% ---------------------------------------

\item \textbf{An uncharged atom of phosphorus has an atomic number of 15 and a mass number of 31. How many valence electrons does it have?}
\begin{enumerate}[label=\Alph*.]
    \item 15
    \item 31
    \item 16
    \item 5
    \item 3
\end{enumerate}
\textbf{Answer:} D \\
\textbf{Explanation:} With 15 electrons total, the shells fill as 2, 8, and 5. The outermost (valence) shell has 5 electrons.

\item \textbf{Nitrogen has an atomic number of 7. What is its typical bonding capacity (valence)?}
\begin{enumerate}[label=\Alph*.]
    \item 1
    \item 2
    \item 3
    \item 4
    \item 5
\end{enumerate}
\textbf{Answer:} C \\
\textbf{Explanation:} Nitrogen has 5 valence electrons. It needs 3 more to complete its octet, so it typically forms 3 covalent bonds.

\item \textbf{Two atoms are interacting. Atom X has an electronegativity of 3.5, and Atom Y has an electronegativity of 0.9. What type of bond will most likely form between them?}
\begin{enumerate}[label=\Alph*.]
    \item Nonpolar covalent
    \item Polar covalent
    \item Ionic
    \item Hydrogen
    \item Van der Waals
\end{enumerate}
\textbf{Answer:} C \\
\textbf{Explanation:} The massive difference in electronegativity (2.6) means Atom X will strip an electron from Atom Y completely, forming an ionic bond.

\item \textbf{A radioactive isotope of carbon ($^{14}C$) decays into nitrogen ($^{14}N$). What specifically occurred in the nucleus?}
\begin{enumerate}[label=\Alph*.]
    \item A proton transformed into a neutron.
    \item A neutron transformed into a proton.
    \item The nucleus captured an electron.
    \item The nucleus lost a full helium atom.
    \item The mass number increased.
\end{enumerate}
\textbf{Answer:} B \\
\textbf{Explanation:} Beta decay in Carbon-14 involves a neutron decaying into a proton (and emitting an electron), which changes the atomic number from 6 to 7.

\item \textbf{As an electron moves from the third shell to the second shell, it:}
\begin{enumerate}[label=\Alph*.]
    \item Absorbs energy.
    \item Releases energy, often as light or heat.
    \item Gains mass.
    \item Becomes a valence electron.
    \item Causes the atom to become an isotope.
\end{enumerate}
\textbf{Answer:} B \\
\textbf{Explanation:} Moving to a shell closer to the nucleus reduces the electron's potential energy, forcing it to release the difference in energy.

\item \textbf{The chemical behavior of an atom is almost entirely determined by its:}
\begin{enumerate}[label=\Alph*.]
    \item Mass number
    \item Number of neutrons
    \item Inner-shell electrons
    \item Valence electrons
    \item Nuclear charge
\end{enumerate}
\textbf{Answer:} D \\
\textbf{Explanation:} Only the outermost electrons (valence electrons) interact with other atoms during chemical reactions.

\item \textbf{Why are noble gases like Helium, Neon, and Argon chemically inert?}
\begin{enumerate}[label=\Alph*.]
    \item They lack neutrons.
    \item Their outer electron shells are completely filled.
    \item They are highly electronegative.
    \item They exist only as radioactive isotopes.
    \item They have no valence shells.
\end{enumerate}
\textbf{Answer:} B \\
\textbf{Explanation:} A full valence shell makes an atom exceptionally stable, meaning it has no tendency to gain, lose, or share electrons.

\item \textbf{Which statement best describes a polar covalent bond?}
\begin{enumerate}[label=\Alph*.]
    \item Electrons are shared equally between atoms.
    \item One atom completely strips an electron from another.
    \item Electrons are shared unequally, creating partial positive and negative charges.
    \item It is a weak interaction between transient dipoles.
    \item It only occurs between ions in a lattice.
\end{enumerate}
\textbf{Answer:} C \\
\textbf{Explanation:} In a polar covalent bond, the more electronegative atom pulls the shared electrons closer, establishing $\delta+$ and $\delta-$ regions.

\item \textbf{Which of the following molecules contains a nonpolar covalent bond?}
\begin{enumerate}[label=\Alph*.]
    \item $H_2O$
    \item $NH_3$
    \item $NaCl$
    \item $O_2$
    \item $HCl$
\end{enumerate}
\textbf{Answer:} D \\
\textbf{Explanation:} In $O_2$, both oxygen atoms have identical electronegativity, so they share the electrons perfectly equally.

\item \textbf{At chemical equilibrium, which of the following is absolutely true?}
\begin{enumerate}[label=\Alph*.]
    \item The concentrations of reactants and products are equal.
    \item The reaction has completely stopped.
    \item All reactants have been converted to products.
    \item The forward and reverse reaction rates are equal.
    \item The energy of the system is at its maximum.
\end{enumerate}
\textbf{Answer:} D \\
\textbf{Explanation:} Equilibrium is dynamic; reactions continue, but the rates of the forward and reverse reactions are identical, stabilizing concentrations.

\item \textbf{A cation is an ion that has:}
\begin{enumerate}[label=\Alph*.]
    \item Gained an electron and is negatively charged.
    \item Lost an electron and is positively charged.
    \item Gained a proton and is positively charged.
    \item Lost a neutron and is lighter.
    \item Shared an electron unequally.
\end{enumerate}
\textbf{Answer:} B \\
\textbf{Explanation:} Cations carry a positive charge because they have lost one or more negatively charged electrons.

\item \textbf{How many total electrons can the $p$ subshell (consisting of three $p$ orbitals) hold?}
\begin{enumerate}[label=\Alph*.]
    \item 2
    \item 4
    \item 6
    \item 8
    \item 10
\end{enumerate}
\textbf{Answer:} C \\
\textbf{Explanation:} Each of the three $p$ orbitals can hold 2 electrons, making a total of 6 electrons.

\item \textbf{What is the primary role of van der Waals interactions in biology?}
\begin{enumerate}[label=\Alph*.]
    \item They form the strong backbone of DNA molecules.
    \item They completely dissociate salts in water.
    \item They allow transient, weak attractions between nonpolar molecules closely pressed together.
    \item They are the primary bond in water molecules.
    \item They cause radioactive decay.
\end{enumerate}
\textbf{Answer:} C \\
\textbf{Explanation:} Fluctuating electron distributions create temporary positive and negative patches that allow molecules (like a gecko's foot pads and a wall) to adhere weakly.

\item \textbf{Which property is considered an emergent property of sodium chloride (NaCl)?}
\begin{enumerate}[label=\Alph*.]
    \item It is a highly reactive, soft metal.
    \item It is a toxic, green gas.
    \item It is an edible, stable crystal that is safe to consume.
    \item It is a liquid at room temperature.
    \item It has the exact same properties as its constituent elements.
\end{enumerate}
\textbf{Answer:} C \\
\textbf{Explanation:} The compound NaCl has completely different properties (edible crystal) than pure Na (explosive metal) and pure Cl (toxic gas).

\item \textbf{An element has 8 protons, 9 neutrons, and 8 electrons. Its mass number and electrical charge are:}
\begin{enumerate}[label=\Alph*.]
    \item 17; neutral
    \item 16; neutral
    \item 17; +1
    \item 16; -1
    \item 8; neutral
\end{enumerate}
\textbf{Answer:} A \\
\textbf{Explanation:} Mass number = protons + neutrons (8+9=17). Because protons equal electrons, the charge is 0 (neutral).

\item \textbf{Isotopes of a given element differ in their number of:}
\begin{enumerate}[label=\Alph*.]
    \item Protons
    \item Neutrons
    \item Electrons
    \item Orbitals
    \item Valence shells
\end{enumerate}
\textbf{Answer:} B \\
\textbf{Explanation:} Isotopes share the same atomic number (protons) but vary in atomic mass due to different numbers of neutrons.

\item \textbf{A double covalent bond, such as that in carbon dioxide ($O=C=O$), involves the sharing of:}
\begin{enumerate}[label=\Alph*.]
    \item One electron
    \item Two electrons
    \item Three electrons
    \item Four electrons
    \item Six electrons
\end{enumerate}
\textbf{Answer:} D \\
\textbf{Explanation:} A single bond shares one pair (2 electrons), so a double bond shares two pairs (4 electrons).

\item \textbf{The shape of a methane molecule ($CH_4$) is a tetrahedron because:}
\begin{enumerate}[label=\Alph*.]
    \item Carbon forms ionic bonds with hydrogen.
    \item The 1s orbital of carbon pushes the hydrogens away.
    \item The hybridization of carbon's one $s$ and three $p$ orbitals creates four teardrop-shaped orbitals spaced evenly apart.
    \item Hydrogen atoms repel each other into a flat square.
    \item Methane is highly polar.
\end{enumerate}
\textbf{Answer:} C \\
\textbf{Explanation:} The $sp^3$ hybridization of carbon's valence shell orbitals results in a tetrahedral geometry with $109.5^\circ$ angles.

\item \textbf{Why does salt dissolve so readily in water?}
\begin{enumerate}[label=\Alph*.]
    \item Water is nonpolar and mixes with nonpolar salts.
    \item The partial charges of water molecules electrostatically surround the individual $Na^+$ and $Cl^-$ ions, breaking the lattice.
    \item Water forms strong covalent bonds with sodium.
    \item Water turns into a gas when touching salt.
    \item Salt lowers the pH of the water drastically.
\end{enumerate}
\textbf{Answer:} B \\
\textbf{Explanation:} Water's polarity creates hydration shells around the charged ions, insulating them from each other and dissolving the crystal.

\item \textbf{Endorphins and morphine have very different chemical formulas but similar biological effects because:}
\begin{enumerate}[label=\Alph*.]
    \item They both dissolve in water.
    \item They share a similar specific three-dimensional shape, allowing them to bind to the same cellular receptors.
    \item They are both radioactive.
    \item They have the exact same number of carbon atoms.
    \item They act as enzymes in the brain.
\end{enumerate}
\textbf{Answer:} B \\
\textbf{Explanation:} Molecular shape dictates function in biology; morphine mimics the shape of endorphins well enough to bind to endorphin receptors.

% ---------------------------------------
% CHAPTER 3: Water (41-60)
% ---------------------------------------

\item \textbf{Which bonds must be broken for water to vaporize into steam?}
\begin{enumerate}[label=\Alph*.]
    \item Polar covalent bonds
    \item Nonpolar covalent bonds
    \item Hydrogen bonds
    \item Ionic bonds
    \item Van der Waals interactions
\end{enumerate}
\textbf{Answer:} C \\
\textbf{Explanation:} Vaporization requires breaking the intermolecular hydrogen bonds holding the liquid together, not the intramolecular covalent bonds holding the $H$ and $O$ atoms together.

\item \textbf{The remarkable high surface tension of water is due primarily to:}
\begin{enumerate}[label=\Alph*.]
    \item The high mass of the oxygen atom.
    \item Hydrophobic interactions at the air-water interface.
    \item The ordered arrangement of hydrogen bonds between water molecules at the surface.
    \item Ionic bonding between dissolved minerals.
    \item The low specific heat of water.
\end{enumerate}
\textbf{Answer:} C \\
\textbf{Explanation:} Molecules at the surface hydrogen-bond to their neighbors and those below them, creating a taut, resilient "film."

\item \textbf{If a lake warms from $1^\circ\text{C}$ to $4^\circ\text{C}$, the water in the lake will:}
\begin{enumerate}[label=\Alph*.]
    \item Freeze solid.
    \item Expand and become less dense.
    \item Contract and become more dense.
    \item Vaporize.
    \item Sublimate.
\end{enumerate}
\textbf{Answer:} C \\
\textbf{Explanation:} Water is unique because it reaches its maximum density at $4^\circ\text{C}$; warming it from $1^\circ\text{C}$ to $4^\circ\text{C}$ causes it to contract.

\item \textbf{Why does sweating effectively cool the human body?}
\begin{enumerate}[label=\Alph*.]
    \item Sweat is naturally colder than body temperature.
    \item Water has a high heat of vaporization, meaning the fastest (hottest) molecules carry a massive amount of body heat away as they turn to gas.
    \item Sweat forms an insulating layer of ice.
    \item Sweat destroys the hydrogen bonds in the skin.
    \item Sweat lowers the pH of the skin.
\end{enumerate}
\textbf{Answer:} B \\
\textbf{Explanation:} Evaporative cooling works because evaporating water requires a significant input of thermal energy to break its hydrogen bonds.

\item \textbf{A solution has a pH of 5. What is the concentration of hydrogen ions $[H^+]$ in this solution?}
\begin{enumerate}[label=\Alph*.]
    \item $10^{-5}$ M
    \item $10^{5}$ M
    \item $10^{-9}$ M
    \item $5$ M
    \item $-5$ M
\end{enumerate}
\textbf{Answer:} A \\
\textbf{Explanation:} By definition, $\text{pH} = -\log[H^+]$, so a pH of 5 means $[H^+] = 10^{-5}$ M.

\item \textbf{How many times more acidic is a solution with a pH of 3 compared to a solution with a pH of 6?}
\begin{enumerate}[label=\Alph*.]
    \item 2 times
    \item 3 times
    \item 30 times
    \item 100 times
    \item 1,000 times
\end{enumerate}
\textbf{Answer:} E \\
\textbf{Explanation:} The pH scale is logarithmic. A difference of 3 units means a $10 \times 10 \times 10 = 1,000$-fold difference in $[H^+]$.

\item \textbf{If $[H^+] = 10^{-4}$ M, what is the $[OH^-]$ of the solution at $25^\circ\text{C}$?}
\begin{enumerate}[label=\Alph*.]
    \item $10^{-10}$ M
    \item $10^{-4}$ M
    \item $10^{-14}$ M
    \item $10^{10}$ M
    \item $10^{-7}$ M
\end{enumerate}
\textbf{Answer:} A \\
\textbf{Explanation:} The ion product of water $[H^+][OH^-] = 10^{-14}$. If $[H^+]$ is $10^{-4}$, $[OH^-]$ must be $10^{-10}$.

\item \textbf{Ocean acidification specifically threatens marine organisms with calcium carbonate shells because:}
\begin{enumerate}[label=\Alph*.]
    \item It directly dissolves the calcium in their bodies.
    \item The increased $CO_2$ blocks sunlight from reaching them.
    \item The extra $H^+$ ions combine with carbonate ($CO_3^{2-}$), reducing the carbonate concentration available for calcification.
    \item It raises the temperature of the water drastically.
    \item It turns the ocean into a strong base.
\end{enumerate}
\textbf{Answer:} C \\
\textbf{Explanation:} Carbonic acid dissociates into $H^+$, which "steals" carbonate ions to form bicarbonate, making it harder for corals to build shells.

\item \textbf{A buffer is a substance that minimizes changes in pH by:}
\begin{enumerate}[label=\Alph*.]
    \item Precipitating out of solution.
    \item Accepting $H^+$ when they are in excess and donating $H^+$ when they have been depleted.
    \item Always keeping the solution at exactly pH 7.0.
    \item Completely neutralizing strong acids into salts.
    \item Boiling off excess carbon dioxide.
\end{enumerate}
\textbf{Answer:} B \\
\textbf{Explanation:} Buffers (usually weak acid-base pairs) reversibly bind and release hydrogen ions to stabilize the solution's pH.

\item \textbf{Which property is necessary for a molecule to be hydrophilic but completely insoluble in water?}
\begin{enumerate}[label=\Alph*.]
    \item It must be a lipid.
    \item It must be small and uncharged.
    \item It must be a massive macromolecule with polar regions, like cellulose.
    \item It must have a high specific heat.
    \item It must be a strong acid.
\end{enumerate}
\textbf{Answer:} C \\
\textbf{Explanation:} Large molecules like cotton (cellulose) can absorb water due to polar regions (hydrophilic) but are too massive to dissolve (insoluble).

\item \textbf{Water's ability to act as a versatile solvent arises from its:}
\begin{enumerate}[label=\Alph*.]
    \item High molecular mass
    \item Polarity and ability to form hydration shells
    \item Low specific heat
    \item Linear molecular shape
    \item Ability to form \textit{cis-trans} isomers
\end{enumerate}
\textbf{Answer:} B \\
\textbf{Explanation:} Water's partial charges allow it to electrostatic surround and separate both ions and polar molecules.

\item \textbf{The specific heat of water is defined as 1 calorie per gram per degree Celsius. Because this is unusually high, water:}
\begin{enumerate}[label=\Alph*.]
    \item Freezes very quickly.
    \item Evaporates at low temperatures.
    \item Resists changes in temperature, providing a stable environment for life.
    \item Cannot be boiled.
    \item Has a very low surface tension.
\end{enumerate}
\textbf{Answer:} C \\
\textbf{Explanation:} Water must absorb or release a massive amount of heat before its temperature changes significantly.

\item \textbf{In a single molecule of water, the oxygen atom is bonded to the hydrogen atoms by:}
\begin{enumerate}[label=\Alph*.]
    \item Hydrogen bonds
    \item Ionic bonds
    \item Polar covalent bonds
    \item Nonpolar covalent bonds
    \item Metallic bonds
\end{enumerate}
\textbf{Answer:} C \\
\textbf{Explanation:} Oxygen's high electronegativity pulls shared electrons toward it, creating polar covalent bonds within the molecule.

\item \textbf{Adhesion is an emergent property of water characterized by:}
\begin{enumerate}[label=\Alph*.]
    \item Water molecules sticking to other water molecules.
    \item Water molecules clinging to another substance, such as the cell walls of plant veins.
    \item The surface tension of a pond.
    \item Water dissolving salt.
    \item The expansion of freezing water.
\end{enumerate}
\textbf{Answer:} B \\
\textbf{Explanation:} Adhesion is the attraction between different substances, which helps water move up against gravity in plants.

\item \textbf{Ice floats on liquid water because:}
\begin{enumerate}[label=\Alph*.]
    \item Ice is heavier than water.
    \item The hydrogen bonds become rigid, forcing molecules further apart and lowering the density.
    \item The covalent bonds expand when cooled.
    \item Ice is hydrophobic and repels liquid water.
    \item Ice loses oxygen atoms when freezing.
\end{enumerate}
\textbf{Answer:} B \\
\textbf{Explanation:} In ice, molecules are locked into a spacious crystalline lattice by stable hydrogen bonds, making ice about 10\% less dense than liquid water.

\item \textbf{A strong acid like HCl differs from a weak acid like carbonic acid ($H_2CO_3$) because:}
\begin{enumerate}[label=\Alph*.]
    \item HCl accepts $H^+$ instead of donating it.
    \item HCl completely dissociates in water, while $H_2CO_3$ reversibly and partially dissociates.
    \item HCl has a pH of 14.
    \item HCl forms hydration shells around lipids.
    \item HCl contains carbon.
\end{enumerate}
\textbf{Answer:} B \\
\textbf{Explanation:} Strong acids undergo complete, irreversible dissociation in water, flooding the solution with $H^+$.

\item \textbf{You have a 1 Molar (1 M) solution of glucose. This means there is:}
\begin{enumerate}[label=\Alph*.]
    \item 1 gram of glucose per liter of solution.
    \item 1 mole of glucose per liter of solution.
    \item 1 mole of glucose per liter of solvent.
    \item 1 gram of glucose per liter of solvent.
    \item 1 molecule of glucose per liter of solution.
\end{enumerate}
\textbf{Answer:} B \\
\textbf{Explanation:} Molarity is defined as the number of moles of a solute per liter of the entire solution.

\item \textbf{What constitutes a single mole of any substance?}
\begin{enumerate}[label=\Alph*.]
    \item Exactly 100 grams.
    \item Exactly 1 liter.
    \item Avogadro's number ($6.02 \times 10^{23}$) of molecules of that substance.
    \item The molecular mass divided by 10.
    \item The total number of protons in the substance.
\end{enumerate}
\textbf{Answer:} C \\
\textbf{Explanation:} A mole is an exact number of objects ($6.02 \times 10^{23}$), providing a bridge between molecular mass in daltons and macro-mass in grams.

\item \textbf{If a base is added to an unbuffered aqueous solution, what immediately happens to the $[H^+]$ and the pH?}
\begin{enumerate}[label=\Alph*.]
    \item $[H^+]$ increases; pH decreases.
    \item $[H^+]$ decreases; pH increases.
    \item Both $[H^+]$ and pH increase.
    \item Both $[H^+]$ and pH decrease.
    \item $[H^+]$ decreases; pH remains constant.
\end{enumerate}
\textbf{Answer:} B \\
\textbf{Explanation:} Bases lower the $H^+$ concentration. Because pH is a negative log scale, a lower $[H^+]$ means a higher pH number.

\item \textbf{Thermal energy differs from temperature in that:}
\begin{enumerate}[label=\Alph*.]
    \item Temperature is total kinetic energy; thermal energy is average kinetic energy.
    \item Thermal energy depends on the total volume of the matter, whereas temperature is an average unaffected by total volume.
    \item Thermal energy is measured in degrees Celsius.
    \item Temperature applies only to gases.
    \item Thermal energy is potential energy, not kinetic.
\end{enumerate}
\textbf{Answer:} B \\
\textbf{Explanation:} A massive, cool ocean has far more thermal energy (total kinetic energy) than a small, hot cup of coffee, even though the coffee has a higher temperature (average kinetic energy).

% ---------------------------------------
% CHAPTER 4: Carbon (61-80)
% ---------------------------------------

\item \textbf{The historical theory of vitalism was fundamentally contradicted by:}
\begin{enumerate}[label=\Alph*.]
    \item The discovery of DNA.
    \item The abiotic synthesis of organic compounds in a laboratory setting.
    \item The observation of cell division under a microscope.
    \item The discovery of isotopes.
    \item The sequencing of the human genome.
\end{enumerate}
\textbf{Answer:} B \\
\textbf{Explanation:} Wöhler's synthesis of urea and Miller's experiment proved that physical and chemical laws, not mystical life forces, govern the formation of organic molecules.

\item \textbf{Which characteristic of carbon atoms is most responsible for the vast molecular diversity of life?}
\begin{enumerate}[label=\Alph*.]
    \item High electronegativity
    \item Ability to form hydrogen bonds
    \item Tetravalence, allowing it to act as an intersection branching in four directions
    \item Its status as a noble gas
    \item Its extremely high mass
\end{enumerate}
\textbf{Answer:} C \\
\textbf{Explanation:} Because carbon can form four covalent bonds, it can build massive, intricate chains, rings, and branched structures.

\item \textbf{Isomers that are mirror images of each other are called:}
\begin{enumerate}[label=\Alph*.]
    \item Structural isomers
    \item Cis-trans isomers
    \item Enantiomers
    \item Radioactive isotopes
    \item Polymers
\end{enumerate}
\textbf{Answer:} C \\
\textbf{Explanation:} Enantiomers differ in spatial arrangement around an asymmetric (chiral) carbon, resulting in non-superimposable left and right-handed mirror images.

\item \textbf{For a carbon atom to be considered asymmetric (chiral), it must be:}
\begin{enumerate}[label=\Alph*.]
    \item Bonded to four identical groups.
    \item Part of a double bond.
    \item Bonded to four completely different atoms or groups of atoms.
    \item Bonded only to hydrogen.
    \item Located at the absolute end of a carbon chain.
\end{enumerate}
\textbf{Answer:} C \\
\textbf{Explanation:} An asymmetric carbon lacks a plane of symmetry because it is bonded to four distinct chemical groups.

\item \textbf{Which requirement is strictly necessary for a molecule to exhibit cis-trans (geometric) isomerism?}
\begin{enumerate}[label=\Alph*.]
    \item An asymmetric carbon.
    \item An inflexible double bond between two carbon atoms.
    \item A ring structure.
    \item A sulfhydryl group.
    \item At least 10 carbon atoms.
\end{enumerate}
\textbf{Answer:} B \\
\textbf{Explanation:} Single bonds allow free rotation, but a double bond locks atoms into place, creating "cis" (same side) and "trans" (opposite sides) arrangements.

\item \textbf{The pharmacological significance of enantiomers is highlighted by the fact that:}
\begin{enumerate}[label=\Alph*.]
    \item They always have the exact same biological effect.
    \item One enantiomer is usually highly radioactive.
    \item Usually, only one enantiomer is biologically active due to the specific 3D shape of cellular receptors.
    \item The body cannot digest any enantiomers.
    \item They always neutralize each other in the bloodstream.
\end{enumerate}
\textbf{Answer:} C \\
\textbf{Explanation:} Receptors are highly shape-specific, just as a right-handed glove only fits a right hand.

\item \textbf{Which functional group is characterized by a carbon atom double-bonded to an oxygen atom?}
\begin{enumerate}[label=\Alph*.]
    \item Hydroxyl
    \item Carbonyl
    \item Amino
    \item Methyl
    \item Sulfhydryl
\end{enumerate}
\textbf{Answer:} B \\
\textbf{Explanation:} The carbonyl group ($>C=O$) is the defining feature of ketones and aldehydes.

\item \textbf{If a carbonyl group is located at the extreme end of a carbon skeleton, the compound is classified as a(n):}
\begin{enumerate}[label=\Alph*.]
    \item Ketone
    \item Aldehyde
    \item Alcohol
    \item Thiol
    \item Amine
\end{enumerate}
\textbf{Answer:} B \\
\textbf{Explanation:} Terminal carbonyls form aldehydes (like formaldehyde), while internal carbonyls form ketones (like acetone).

\item \textbf{Which functional group acts as an acid by easily donating a hydrogen ion ($H^+$)?}
\begin{enumerate}[label=\Alph*.]
    \item Hydroxyl ($-OH$)
    \item Carboxyl ($-COOH$)
    \item Amino ($-NH_2$)
    \item Sulfhydryl ($-SH$)
    \item Methyl ($-CH_3$)
\end{enumerate}
\textbf{Answer:} B \\
\textbf{Explanation:} The two highly electronegative oxygen atoms in the carboxyl group pull electrons away from hydrogen so strongly that the $H^+$ dissociates.

\item \textbf{The amino functional group acts as a biological base because it:}
\begin{enumerate}[label=\Alph*.]
    \item Donates $OH^-$ ions to the solution.
    \item Absorbs $H^+$ from the surrounding solution, becoming $-NH_3^+$.
    \item Completely dissociates in water.
    \item Acts as a buffer against all pH changes.
    \item Releases pure nitrogen gas.
\end{enumerate}
\textbf{Answer:} B \\
\textbf{Explanation:} By picking up an $H^+$ from the solution, the amino group decreases the $[H^+]$ concentration, fulfilling the definition of a base.

\item \textbf{A molecule with the chemical formula $C_3H_8$ is a(n):}
\begin{enumerate}[label=\Alph*.]
    \item Carbohydrate
    \item Amino acid
    \item Hydrocarbon
    \item Nucleic acid
    \item Steroid
\end{enumerate}
\textbf{Answer:} C \\
\textbf{Explanation:} It contains exclusively carbon and hydrogen, making it a pure hydrocarbon (propane).

\item \textbf{Which two functional groups are always present in every amino acid?}
\begin{enumerate}[label=\Alph*.]
    \item Carbonyl and Hydroxyl
    \item Amino and Carboxyl
    \item Sulfhydryl and Phosphate
    \item Methyl and Carbonyl
    \item Amino and Sulfhydryl
\end{enumerate}
\textbf{Answer:} B \\
\textbf{Explanation:} The very name "amino acid" refers to the basic amino group and the acidic carboxyl group attached to the alpha carbon.

\item \textbf{Which functional group is highly unreactive but serves as a crucial biological tag, particularly in regulating gene expression?}
\begin{enumerate}[label=\Alph*.]
    \item Phosphate
    \item Carboxyl
    \item Methyl
    \item Sulfhydryl
    \item Hydroxyl
\end{enumerate}
\textbf{Answer:} C \\
\textbf{Explanation:} The nonpolar methyl group ($-CH_3$) doesn't participate in standard chemical reactions but heavily influences the shape and function of DNA and sex hormones.

\item \textbf{The sulfhydryl group ($-SH$) is vital for protein structure because:}
\begin{enumerate}[label=\Alph*.]
    \item It provides the energy for peptide bond formation.
    \item Two sulfhydryl groups can react to form a strong, covalent "cross-link" (disulfide bridge).
    \item It makes the protein highly soluble in lipids.
    \item It acts as the primary buffer in blood.
    \item It forms the genetic code.
\end{enumerate}
\textbf{Answer:} B \\
\textbf{Explanation:} Cysteine amino acids use their sulfhydryl groups to rivet folded proteins into their permanent 3D tertiary structures.

\item \textbf{ATP stores potential energy that can be released for cellular work. Which functional group is primarily responsible for this energy transfer?}
\begin{enumerate}[label=\Alph*.]
    \item Sulfhydryl
    \item Methyl
    \item Phosphate
    \item Amino
    \item Carbonyl
\end{enumerate}
\textbf{Answer:} C \\
\textbf{Explanation:} Adenosine triphosphate (ATP) contains three clustered phosphate groups. The hydrolysis of the terminal phosphate releases a large amount of energy.

\item \textbf{Differences in the skeletal structure between pentane ($C_5H_{12}$) and isopentane ($C_5H_{12}$) make them:}
\begin{enumerate}[label=\Alph*.]
    \item Enantiomers
    \item Structural isomers
    \item Cis-trans isomers
    \item Radioactive isotopes
    \item Unrelated molecules
\end{enumerate}
\textbf{Answer:} B \\
\textbf{Explanation:} They have the same formula but differ in the covalent arrangement of their carbon skeletons (straight chain vs. branched).

\item \textbf{Which statement regarding hydrocarbons is true?}
\begin{enumerate}[label=\Alph*.]
    \item They are highly hydrophilic.
    \item They readily form hydrogen bonds with water.
    \item They consist of highly polar covalent bonds.
    \item They are hydrophobic and store a large amount of energy.
    \item They are all solid at room temperature.
\end{enumerate}
\textbf{Answer:} D \\
\textbf{Explanation:} The C-H bond is nonpolar, making hydrocarbons hydrophobic. Breaking these bonds releases substantial energy (as seen in gasoline and fat tails).

\item \textbf{In Stanley Miller's experiment, what was the significance of finding amino acids in the collection flask?}
\begin{enumerate}[label=\Alph*.]
    \item It proved that life currently exists on other planets.
    \item It showed that organic molecules can be synthesized abiotically under conditions mimicking early Earth.
    \item It proved that DNA is older than proteins.
    \item It demonstrated that oxygen was plentiful on early Earth.
    \item It showed that vitalism is the only valid theory of life.
\end{enumerate}
\textbf{Answer:} B \\
\textbf{Explanation:} This was the defining blow to vitalism, proving organic compounds could arise spontaneously through natural chemical and physical processes.

\item \textbf{A compound ending in the suffix "-ol" (e.g., ethanol) most likely contains which functional group?}
\begin{enumerate}[label=\Alph*.]
    \item Carboxyl
    \item Amino
    \item Sulfhydryl
    \item Hydroxyl
    \item Methyl
\end{enumerate}
\textbf{Answer:} D \\
\textbf{Explanation:} Compounds with hydroxyl groups are classified as alcohols, which are typically named with an "-ol" ending.

\item \textbf{The primary difference between estradiol and testosterone, two vastly different hormones, is:}
\begin{enumerate}[label=\Alph*.]
    \item They have completely different four-ring carbon skeletons.
    \item Estradiol is a protein, and testosterone is a lipid.
    \item They share the same carbon skeleton but differ in the specific chemical groups attached to those rings.
    \item Testosterone lacks carbon atoms.
    \item Estradiol is highly acidic.
\end{enumerate}
\textbf{Answer:} C \\
\textbf{Explanation:} Both are steroids derived from cholesterol. Their drastic biological differences stem entirely from slight variations in their attached functional groups.

% ---------------------------------------
% CHAPTER 5: Macromolecules (81-100)
% ---------------------------------------

\item \textbf{Which of the following processes requires a net input of energy and the removal of a water molecule?}
\begin{enumerate}[label=\Alph*.]
    \item Dehydration synthesis (condensation reaction)
    \item Hydrolysis
    \item Denaturation
    \item Renaturation
    \item Isomerization
\end{enumerate}
\textbf{Answer:} A \\
\textbf{Explanation:} Dehydration synthesis links monomers together, a process that inherently requires energy and the extraction of $H_2O$.

\item \textbf{A molecule of dietary fiber (cellulose) cannot be digested by humans because:}
\begin{enumerate}[label=\Alph*.]
    \item It is a lipid, not a carbohydrate.
    \item Humans lack the specific enzymes necessary to hydrolyze its $\beta$ (beta) glycosidic linkages.
    \item It contains toxic nitrogenous bases.
    \item It is made of completely rigid disulfide bridges.
    \item Humans lack an acidic stomach.
\end{enumerate}
\textbf{Answer:} B \\
\textbf{Explanation:} Human amylase enzymes can only break the $\alpha$ (alpha) linkages found in starch and glycogen, not the straight $\beta$ linkages of cellulose.

\item \textbf{Which structural feature allows parallel cellulose molecules to form incredibly tough microfibrils in plant cell walls?}
\begin{enumerate}[label=\Alph*.]
    \item Extensive branching of the carbon skeleton.
    \item Covalent cross-links between strands.
    \item Hydrogen bonding between the hydroxyl groups of adjacent, unbranched parallel strands.
    \item High lipid content.
    \item Double-helical twisting.
\end{enumerate}
\textbf{Answer:} C \\
\textbf{Explanation:} Because cellulose is straight (unbranched), parallel molecules align closely and lock together via numerous hydrogen bonds.

\item \textbf{What distinguishes saturated fats from unsaturated fats at the molecular level?}
\begin{enumerate}[label=\Alph*.]
    \item Saturated fats have long hydrocarbon tails; unsaturated fats do not.
    \item Saturated fats contain only single carbon-carbon bonds, maximizing hydrogen content, while unsaturated fats contain one or more \textit{cis} double bonds.
    \item Saturated fats are highly hydrophilic.
    \item Unsaturated fats contain multiple phosphate groups.
    \item Saturated fats are exclusively derived from plants.
\end{enumerate}
\textbf{Answer:} B \\
\textbf{Explanation:} The lack of double bonds allows saturated fat chains to pack tightly into solids. The \textit{cis} double bonds in unsaturated fats create kinks, keeping them liquid.

\item \textbf{A trans fat is created through the process of partial hydrogenation. This process:}
\begin{enumerate}[label=\Alph*.]
    \item Converts saturated fats into liquids.
    \item Removes hydrogen atoms to create triple bonds.
    \item Synthetically adds hydrogen to \textit{cis}-unsaturated fats, inadvertently creating unnatural \textit{trans} double bonds that pack tightly.
    \item Converts lipids into proteins.
    \item Turns fatty acids into steroids.
\end{enumerate}
\textbf{Answer:} C \\
\textbf{Explanation:} Hydrogenation is designed to solidify liquid oils, but the resulting \textit{trans} geometric isomers are strongly linked to cardiovascular disease.

\item \textbf{Phospholipids spontaneously form a bilayer in aqueous environments because:}
\begin{enumerate}[label=\Alph*.]
    \item They are entirely hydrophobic and clump together.
    \item They are joined by strong peptide bonds.
    \item They are amphipathic; their hydrophilic heads face the water, while their hydrophobic tails cluster together away from the water.
    \item The water aggressively forms covalent bonds with the tails.
    \item They are naturally magnetic.
\end{enumerate}
\textbf{Answer:} C \\
\textbf{Explanation:} This self-assembly is driven by hydrophobic exclusion, forming the fundamental architecture of all cell membranes.

\item \textbf{The overall three-dimensional shape of a single, completely folded polypeptide chain is defined as its:}
\begin{enumerate}[label=\Alph*.]
    \item Primary structure
    \item Secondary structure
    \item Tertiary structure
    \item Quaternary structure
    \item Pentameric structure
\end{enumerate}
\textbf{Answer:} C \\
\textbf{Explanation:} Tertiary structure is the overall 3D shape resulting from interactions between the diverse R groups (side chains) of the amino acids.

\item \textbf{Which type of bond is strictly responsible for maintaining the secondary structure (like an $\alpha$ helix) of a protein?}
\begin{enumerate}[label=\Alph*.]
    \item Disulfide bridges between cysteine R groups.
    \item Hydrogen bonds between atoms of the repeating polypeptide backbone.
    \item Ionic bonds between acidic and basic R groups.
    \item Peptide bonds.
    \item Phosphodiester linkages.
\end{enumerate}
\textbf{Answer:} B \\
\textbf{Explanation:} Secondary structure relies exclusively on hydrogen bonding between the partially charged oxygen and nitrogen/hydrogen atoms of the uniform backbone, independent of the R groups.

\item \textbf{The quaternary structure of a protein requires:}
\begin{enumerate}[label=\Alph*.]
    \item At least one disulfide bridge.
    \item A specific sequence of strictly polar amino acids.
    \item The aggregation of two or more distinct polypeptide subunits into one functional macromolecule.
    \item Binding to a strand of DNA.
    \item Exposure to extreme heat.
\end{enumerate}
\textbf{Answer:} C \\
\textbf{Explanation:} Proteins like hemoglobin or collagen only achieve their final functional form when multiple separate polypeptide chains assemble together.

\item \textbf{Which of the following factors generally does NOT cause protein denaturation?}
\begin{enumerate}[label=\Alph*.]
    \item Extreme changes in pH.
    \item High temperature.
    \item High salt concentration.
    \item Enzymatic hydrolysis into individual amino acids.
    \item Transfer to a nonpolar solvent.
\end{enumerate}
\textbf{Answer:} D \\
\textbf{Explanation:} Denaturation is the unfolding of the 3D shape while the backbone remains intact. Hydrolysis actually breaks the covalent peptide bonds, destroying the primary structure itself.

\item \textbf{Chitin is to arthropods as \_\_\_\_\_\_\_ is to plants.}
\begin{enumerate}[label=\Alph*.]
    \item Starch
    \item Cellulose
    \item Glycogen
    \item Cholesterol
    \item Peptidoglycan
\end{enumerate}
\textbf{Answer:} B \\
\textbf{Explanation:} Both are structural polysaccharides used to build tough external walls (exoskeletons in arthropods, cell walls in plants).

\item \textbf{The bond formed between two nucleotides in a polynucleotide strand is a:}
\begin{enumerate}[label=\Alph*.]
    \item Glycosidic linkage
    \item Peptide bond
    \item Phosphodiester linkage
    \item Ester linkage
    \item Hydrogen bond
\end{enumerate}
\textbf{Answer:} C \\
\textbf{Explanation:} A phosphate group covalently links the sugars of two adjacent nucleotides, forming the sugar-phosphate backbone.

\item \textbf{Which component is present in RNA but completely absent in DNA?}
\begin{enumerate}[label=\Alph*.]
    \item Cytosine
    \item A phosphate group
    \item Deoxyribose
    \item Uracil
    \item Purines
\end{enumerate}
\textbf{Answer:} D \\
\textbf{Explanation:} RNA substitutes the pyrimidine Uracil (U) in place of DNA's Thymine (T). Additionally, RNA uses the sugar ribose instead of deoxyribose.

\item \textbf{In a DNA double helix, if one strand runs 5' to 3', the complementary strand must run:}
\begin{enumerate}[label=\Alph*.]
    \item 5' to 3' parallel.
    \item 3' to 5' antiparallel.
    \item 5' to 5'.
    \item 3' to 3'.
    \item In a random orientation.
\end{enumerate}
\textbf{Answer:} B \\
\textbf{Explanation:} The two strands of the double helix run in opposite directions (antiparallel) to allow the nitrogenous bases to align and hydrogen-bond correctly.

\item \textbf{A protein sequence is highly conserved across humans, chimpanzees, and mice, but very different in fish. This supports the concept of:}
\begin{enumerate}[label=\Alph*.]
    \item Vitalism
    \item Convergent evolution
    \item Bioinformatics
    \item Descent from a common ancestor (Molecular genealogy)
    \item Denaturation
\end{enumerate}
\textbf{Answer:} D \\
\textbf{Explanation:} Species that diverged more recently share a greater proportion of their DNA and protein sequences, providing strong molecular evidence for evolution.

\item \textbf{Which lipid class consists of a carbon skeleton characterized by four fused rings?}
\begin{enumerate}[label=\Alph*.]
    \item Waxes
    \item Triglycerides
    \item Phospholipids
    \item Steroids
    \item Trans fats
\end{enumerate}
\textbf{Answer:} D \\
\textbf{Explanation:} All steroids, including cholesterol and sex hormones, share this fundamental four-ring structure.

\item \textbf{If a mutation substitutes a highly hydrophobic amino acid for a highly hydrophilic one on the exterior surface of a protein in an aqueous solution, what is the most likely outcome?}
\begin{enumerate}[label=\Alph*.]
    \item The protein will function exactly the same.
    \item The protein will likely misfold as the hydrophobic amino acid is driven toward the core by the surrounding water.
    \item The primary structure will reverse direction.
    \item The entire protein will immediately hydrolyze.
    \item The protein will become a nucleic acid.
\end{enumerate}
\textbf{Answer:} B \\
\textbf{Explanation:} Hydrophobic exclusion forces nonpolar side chains to the interior of the protein. Placing one on the surface disrupts normal folding.

\item \textbf{Which of the following nitrogenous bases is a purine?}
\begin{enumerate}[label=\Alph*.]
    \item Cytosine
    \item Thymine
    \item Uracil
    \item Guanine
    \item Pyridoxine
\end{enumerate}
\textbf{Answer:} D \\
\textbf{Explanation:} Adenine and Guanine are purines (larger, double-ring structures). Cytosine, Thymine, and Uracil are pyrimidines.

\item \textbf{Chaperonins assist with protein folding by:}
\begin{enumerate}[label=\Alph*.]
    \item Directly forming peptide bonds between amino acids.
    \item Providing an isolated, hydrophilic environment shielded from the "bad influences" of the cytoplasm.
    \item Enzymatically breaking down misfolded proteins.
    \item Providing the genetic code to the ribosome.
    \item Adding methyl groups to the finished protein.
\end{enumerate}
\textbf{Answer:} B \\
\textbf{Explanation:} Chaperonins do not dictate the fold; they act as a "shelter" allowing the polypeptide to fold spontaneously without unwanted interactions with other cellular molecules.

\item \textbf{Which of the following represents the correct base-pairing rules within a DNA double helix?}
\begin{enumerate}[label=\Alph*.]
    \item Adenine pairs with Guanine; Cytosine pairs with Thymine.
    \item Adenine pairs with Uracil; Guanine pairs with Cytosine.
    \item Adenine pairs with Cytosine; Guanine pairs with Thymine.
    \item Adenine pairs with Thymine; Guanine pairs with Cytosine.
    \item Purines strictly pair with other Purines.
\end{enumerate}
\textbf{Answer:} D \\
\textbf{Explanation:} Strict complementary base pairing ensures that a purine always pairs with a pyrimidine: A with T (2 hydrogen bonds), and G with C (3 hydrogen bonds).

\end{enumerate}

\end{document}
